\documentclass{article}

\usepackage{fullpage}
\usepackage{url}
\usepackage{graphicx}
\usepackage{amsmath}
\usepackage{physics}
\usepackage{mathtools}
\usepackage{bbold}

\usepackage{lmodern}
\usepackage[T1]{fontenc}
\usepackage[fontsize=13.5pt]{fontsize}

\DeclareMathOperator{\erfi}{erfi}
\newcommand{\R}{\ensuremath{\mathbb{R}}}
\newcommand{\C}{\ensuremath{\mathbb{C}}}

\title{Complex Systems Excercise 1}
\author{Idan Stark}
\date{\today}

\begin{document}

\maketitle

\section{ Dimensional analysis - Flow in porous media}

In this section we consider the 1-dimensional flow of ground water through porous rock above an impenetrable bed. The height of the water above the impenetrable bed is given by $h(x, t)$. The head of the water is the pressure generated by the water height $h$ and is given by $H = \rho gh$. It satisfies the equation
\begin{equation}
    \partial_t H = \kappa \partial_x^2 H^2
\end{equation}
where the the coefficient $\kappa$ depends on the porosity and permeability of the rock and the viscosity of the fluid, but is constant in time. Using the initial uniform background head of the groundwater $H_i = H(x, t = 0)$ and that at $t = 0$ a very large volume of water $I$ is deposited in a narrow region of width $2l$.

\subsection{Relevant parameters}
In the equation above, there are four relevant variables: $t$ in the time derivative, $x$ in the spatial derivative, $H$ is the head, with units of pressure, and $\kappa$, the units of which can be extracted from the equation:
\begin{equation}
\begin{gathered}
    [t] = T
    \qquad
    [x] = L
    \qquad
    [H] = [H_i] = \frac{M}{LT^2}
    \\
    [\kappa] = \frac{[\partial_t H]}{[\partial_x^2 H^2]} = \frac{M/LT^3}{M^2/L^4T^4} = \frac{L^3T}{M}
\end{gathered}
\end{equation}
Additionaly to the variables in the equation, we have a large volume of water deposited $I$ with units $[I] = L^2$. The narrow region in which they are deposited $2l$, can be neglected. In total we have six variables and three scales, resulting in 3 dimensionless variables.
\newline
For the basis of these quantitied, we can use $\kappa$ and $t$, but we need a third variable. Since $H$ is the only other variable with mass dependence, we should select it, but it is the final result we wish to obtain. Instead, let us look at the total $H$, which is a constant, since
\begin{equation}
    W = \int H dx = \rho g \int h dx = \rho g I
\end{equation}
\begin{equation}
    [W] = [H]L = \frac{M}{T^2}
\end{equation}
So, instead of $I$ we use $W$ our three base variables are $\{\kappa, W, t\}$. Using these variables to get the units of the three other variables:
\begin{equation}
    [\kappa]^\alpha[W]^\beta [t]^\gamma = M^{\beta - \alpha}L^{3\alpha}T^{\alpha - 2\beta + \gamma}
\end{equation}
\begin{equation}
    [H] = [H_i] = M^1 L^{-1} T^{-2} \rightarrow
    \alpha = -\frac{1}{3}, \beta = \frac{2}{3}, \gamma = -\frac{1}{3}
\end{equation}
\begin{equation}
    [x] = L \rightarrow \alpha = \beta = \gamma = \frac{1}{3}
\end{equation}
\begin{equation}
    \boxed{
        \pi_1 = H\left(\frac{\kappa t}{W^2}\right)^{1/3}
        \qquad
        \pi_2 = \frac{x}{(\kappa W t)^{1/3}}
        \qquad
        \pi_3 = H_i\left(\frac{\kappa t}{W^2}\right)^{1/3}
    }
\end{equation}
\subsection{Self similar form}
These dimensionless variables must rely on each other:
\begin{equation}
    \pi_1 = f(\pi_2, \pi_3)
\end{equation}
But $\pi_3$ only describes the relations of the starting conditions. After a long enough time, we can assume these conditions no longer matter and thus we can ignore this variable:
\begin{equation}
    \pi_1 = f(\pi_2)
\end{equation}
This gives us:
\begin{equation}
    H = \left(\frac{W^2}{\kappa t}\right)^{1/3} f(\xi)
\end{equation}
Where $\xi = \pi_2 = \frac{x}{(\kappa W t)^{1/3}}$. It will be easy to name the cubic root
\begin{equation}
    L = (\kappa W t)^{1/3}
\end{equation}
to get:
\begin{equation}
    H = \frac{W}{L} f(\xi)
    \qquad
    \xi = \frac{x}{L}
\end{equation}
Now, we can assign this back into the equation to get:
\begin{equation}
\begin{gathered}
    \partial_t H = -\frac{1}{3t}\frac{W}{L} \left(f(\xi) + \xi f'(\xi)\right)
    \\
    \partial_x^2 H^2 = 
    \frac{W^2}{L^4} \partial_\xi^2 f^2(\xi) =
    2\frac{W^2}{L^4} \left[f(\xi)f''(\xi) + f'^2(\xi)\right]
\end{gathered}
\end{equation}
Which gives:
\begin{equation}
    -\frac{1}{3t}\frac{W}{L} \left(f(\xi) + \xi f'(\xi)\right) =  2\kappa\frac{W^2}{L^4} \left[f(\xi)f''(\xi) + f'^2(\xi)\right]
\end{equation}
Writing in the RHS denominator $L^4 = L\kappa W t$ gives
\begin{equation}
    -\frac{1}{3t}\frac{W}{L} \left(f(\xi) + \xi f'(\xi)\right) =  2\kappa\frac{W^2}{L\kappa At} \left[f(\xi)f''(\xi) + f'^2(\xi)\right]
\end{equation}
\begin{equation}
    \boxed{6\left[f(\xi)f''(\xi) + f'^2(\xi)\right] + \left[f(\xi) + \xi f'(\xi)\right] = 0}
\end{equation}

\subsection{The Solution}
It can be seen that each of the two terms in the equation is a full derivative:
\begin{equation}
    \left[f(\xi)f''(\xi) + f'^2(\xi)\right] = \frac{d}{d\xi}\left[f(\xi)f'(\xi)\right]
    \qquad
    \left[f(\xi) + \xi f'(\xi)\right] = \frac{d}{d\xi}\left[\xi f(\xi)\right]
\end{equation}
Integrating this gives:
\begin{equation}
    6f(\xi)f'(\xi) + \xi f(\xi) = C
\end{equation}
To simplify this, let us require that the solution $H$ is symmetric around $x = 0$, which is expected from the symmetry of the problem. This means that at any given time, $f(\xi)$ is even, and so $f'(0) = 0$. Plugging this in the above equation gives:
\begin{equation}
    C = 6f(0)f'(0) + 0\cdot f(0) = 0
\end{equation}
Which reduces it to a homogenous equation, with the simple solution
\begin{equation}
    f(\xi) = C - \frac{\xi^2}{12}
\end{equation}
Now, we can again use
\begin{equation}
    I = \frac{1}{\rho g} \int H dx = \frac{I}{w} \frac{w}{L} L \int  f(\xi) d\xi \rightarrow \int f(\xi) d\xi = 1
\end{equation}
To get the actual amount of water, we must integrate only over the positive area, bewteen $\xi_0 = \pm \sqrt{12C}$:
\begin{equation}
    1 = \int_{-\sqrt{12C}}^{\sqrt{12C}} C - \frac{\xi^2}{12} = C\xi - \frac{\xi^3}{36}\Big|_{-\sqrt{12C}}^{\sqrt{12C}} = \sqrt{12C}\left(2C - \frac{24C}{36}\right) = \frac{4\sqrt{12}}{3}C^{3/2}
\end{equation}
\begin{equation}
    C = \frac{3^{1/3}}{4}
\end{equation}
And so the final solution gives:
\begin{equation}
    \boxed{H = \left(\frac{\left(\rho g I\right)^2}{\kappa t}\right)^{1/3} \max\left\{\frac{3^{1/3}}{4} - \frac{x^2}{12(\kappa \rho g I t)^{2/3}}, 0\right\}}
\end{equation}

\newpage
\section{Dimensional analysis - Klomogorov scaling}
\subsection{The scaling of Klomogorov length, time and velocity}
Since the Klomogorov scales are dependent only on the energy dissipation rate per unit mass and the kinematic viscosity, we can write it as
\begin{equation}
    l_k = \nu^\alpha\epsilon^\beta
    \qquad
    \tau_k = \nu^\gamma \epsilon^\delta
    \qquad
    v_k = \nu^{\alpha - \gamma}\epsilon^{\beta - \delta}
\end{equation}
Now, using their dimensions
\begin{equation}
    [\nu] = L^{2}T^{-1}
    \qquad
    [\epsilon] = L^2T^{-3}
\end{equation}
We get:
\begin{equation}
    \begin{cases}
        2\alpha + 2\beta = 1
        \\
        -\alpha - 3\beta = 0
    \end{cases}
    \rightarrow
    \begin{matrix}
        \alpha = 3/4 \\ \beta = -1/4
    \end{matrix}
\end{equation}
\begin{equation}
    \begin{cases}
        2\gamma + 2\delta = 0
        \\
        -\gamma - 3\delta = 1
    \end{cases}
    \rightarrow
    \begin{matrix}
        \gamma = 1/2 \\ \delta = -1/2
    \end{matrix}
\end{equation}
Which means:
\begin{equation}
    l_k = \left(\frac{\nu^3}{\epsilon}\right)^{1/4}
    \qquad
    \tau_k = \sqrt{\frac{\nu}{\epsilon}}
    \qquad
    v_k = \left(\nu\epsilon\right)^{1/4}
\end{equation}
\subsection{Reynold's number}
This gives the Reynold's number is
\begin{equation}
    Re = \frac{l_k v_k}{\nu} = \frac{\nu^{3/4}\epsilon^{-1/4}\nu^{1/4}\epsilon^{1/4}}{\nu} = 1
\end{equation}

\subsection{The Energy spectra}
The energy spectra per unit mass $E(k)$ has units of energy per unit mass $k$, and since $[k] = L^{-1}$:
\begin{equation}
    [E(k)] = L^{3}T^{-2}
\end{equation}
Using dimensional analysis, it can be written based on $k$ and $\epsilon$:
\begin{equation}
    E(k) = \epsilon^\alpha k^\beta
\end{equation}
\begin{equation}
    L^{3}T^{-2} = \left(L^2T^{-3}\right)^\alpha \left(L^{-1}\right)^{\beta}
\end{equation}
\begin{equation}
    \begin{cases}
        3 = 2\alpha - \beta \\ -2 = -3\alpha
    \end{cases}
    \rightarrow
    \begin{matrix}
        \alpha = 2/3 \\ \beta = -5/3
    \end{matrix}
\end{equation}
And so
\begin{equation}
    E(k) \propto \epsilon^{2/3} k^{-5/3}
\end{equation}

\newpage
\section{Rank four isotropic tensor}
\subsection{Symmetry in the first two indicies}
Let us take a four rank tensor $T^{ijkl}$ that is symmetric in the first two indicies:
\begin{equation}
    T^{ijkl} = T^{jikl}
\end{equation}
And let us assume it is independent of observer. Specifically, let us assume the tensor is independent of observer:
\begin{equation}
    T^{ijkl} = T^{mnpq}\Lambda^i_m\Lambda^j_n\Lambda^k_p\Lambda^l_q
\end{equation}
For all 3D rotations $\Lambda$. Let us take, for instance, the $\pi/2$ rotations around each axis:
\begin{equation}
    \Lambda_x = \begin{pmatrix}
        1 & 0 & 0 \\ 0 & 0 & 1 \\ 0 & -1 & 0
    \end{pmatrix}
    \qquad
    \Lambda_y = \begin{pmatrix}
        0 & 0 & -1 \\ 0 & 1 & 0 \\ 1 & 0 & 0
    \end{pmatrix}
    \qquad
    \Lambda_z = \begin{pmatrix}
        0 & 1 & 0 \\ -1 & 0 & 0 \\ 0 & 0 & 1
    \end{pmatrix}
\end{equation}
In each rotation, since there is only one non-zero element in each row and column of the matrix, only one term survives in the right hand side for each index combination on the left hand side. If the index matches the rotation axis, it doesn't change. The other indices switch places, changing sign according to the direction of rotation.
\newline
Let us look what happens when there is an odd number of a certain index. In that case, we can use one of the rotations twice to see that it cancels out\footnote{
    Of course this can also be seen using the $\pi$ rotation matrix:
    \begin{equation}
        \Lambda_x(\pi) = \begin{pmatrix}
            1 & 0 & 0 \\ 0 & -1 & 0  \\ 0 & 0 & -1
        \end{pmatrix}
    \end{equation}
Giving:
\begin{equation}
    T^{1222} = \left\{\text{ x rotation }\right\} = -T^{1222}
\end{equation}
But I wanted to use the generic axes-mixing matrices.
}:
\begin{equation}
    T^{1222} = \left\{\text{ x rotation }\right\} = -T^{1333} = \left\{\text{ x rotation }\right\} = -T^{1222} = 0
\end{equation}
So the only possible observer independent tensors are the one where the indices come in pairs. There are three such combination:
\begin{equation}
    T^{ijkl} = \alpha \delta^{ij}\delta^{kl} + \beta \delta^{ik}\delta^{jl} + \gamma \delta^{il}\delta^{jk}
\end{equation}
Now, let us enforce the symmetry in the first two indices:
\begin{equation}
    T^{ijkl} = T^{jikl}
\end{equation}
\begin{equation}
    \alpha \delta^{ij}\delta^{kl} + \beta \delta^{ik}\delta^{jl} + \gamma \delta^{il}\delta^{jk} =
    \alpha \delta^{ji}\delta^{kl} + \beta \delta^{jk}\delta^{il} + \gamma \delta^{jl}\delta^{ik}
\end{equation}
Which results in
\begin{equation}
    \beta = \gamma
\end{equation}
Remaing these constants give:
\begin{equation}
    \boxed{T^{ijkl} = \lambda \delta^{ij}\delta^{kl} + \mu(\delta^{ik}\delta^{jl} + \delta^{jk}\delta^{il})}
\end{equation}
\subsection{Without the symmetry}
As discuseed in the last section, without the symmetry, we get the more generic tensor:
\begin{equation}
    T^{ijkl} = \alpha \delta^{ij}\delta^{kl} + \beta \delta^{ik}\delta^{jl} + \gamma \delta^{il}\delta^{jk}
\end{equation}
with
\begin{equation}
    \boxed{\text{3 degrees of freedom}}
\end{equation}
\subsection{Symmetry in the second pair}
Requiring a symmetry in the second pair of indices:
\begin{equation}
    T^{ijkl} = T^{ijlk}
\end{equation}
Actually doesn't change the tensor at all, since it already has that symmetry. This means it still has
\begin{equation}
    \boxed{\text{2 degrees of freedom}}
\end{equation}

\newpage
\section{Stokes second problem}
In this section, instead of a constant velocity like in Stokes first problem $\boldsymbol{v}(y = 0) = v_0 \boldsymbol{\hat{x}}$, we have an oscillating velocity:
\begin{equation}
    \boldsymbol{v}(y = 0) = v_0 \cos(\omega t) \boldsymbol{\hat{x}}
\end{equation}
Similarly to Stokes first problem, since $\boldsymbol{\nabla}\cdot\boldsymbol{v} = 0$ and from the x-z symmetry, the entire problem reduces to the change along the $\boldsymbol{\hat{y}}$ axis of the $\boldsymbol{v} = v_0\boldsymbol{\hat{x}}$, with constant pressure $p = p_0$. Thus, the Navier-Stokes equations reduce to:
\begin{equation}
    \partial_t v_x = \nu \partial_y^2 v_x
\end{equation}
In addition to the variables showing up in Stokes first problem, we have an additional variable $\omega$ which gives and additional relevant timescale. In total, the four dimensionless variables are:
\begin{equation}
    \pi_1 = \frac{v}{v_0}
    \qquad
    \pi_2 = \sqrt{\frac{\omega}{\nu}}y
    \qquad
    \pi_3 = \omega t
    \qquad
    \pi_4 = \frac{y}{v_0 t}
\end{equation}
Using the same logic as in the first problem, we can assume that the scale of the initial conditions doesn't matter, and thus the total dependence of the variables doesn't include $\pi_4$:
\begin{equation}
    \pi_1 = f(\pi_2, \pi_3) \rightarrow v = v_0 f(\xi, \omega t), \quad \xi = \pi_2
\end{equation}
We can further assume seperability, since we expect $\omega$ to set the timescale of the problem:
\begin{equation}
    v = v_0 g(\omega t) f(\xi)
\end{equation}
Now, using the boundary condition at $y = 0$, we get:
\begin{equation}
    v_0 g(\omega t) f(0) = v_0 \cos(\omega t) \rightarrow g(\omega t) = \frac{1}{f(0)}\cos(\omega t)
\end{equation}
It would be useful to rewrite $\cos(\omega t)$ in its complex form:
\begin{equation}
    v = v_0 e^{i \omega t} \frac{f(\xi)}{f(0)}
\end{equation}
Plugging this back into the equation, multiplying by $f(0)$ and dividing by $v_0$ yields:
\begin{equation}
    \partial_t e^{i \omega t} f(\xi) = \nu \partial_y^2 e^{i \omega t} f(\xi)
\end{equation}
\begin{equation}
    i\omega f(\xi)e^{i \omega t} = e^{i \omega t} \nu \frac{\omega}{\nu} f''(\xi)
\end{equation}
\begin{equation}
    f''(\xi) = if(\xi)
\end{equation}
\begin{equation}
    f(\xi) = e^{\pm\sqrt{i}\xi}
\end{equation}
Now, since
\begin{equation}
    \pm \sqrt{i} = \pm\frac{1 + i}{\sqrt{2}}
\end{equation}
We get:
\begin{equation}
    v(y, t) = v_0 \exp[\pm \sqrt{\frac{\omega}{2\nu}}y +i\left(
        \omega t \pm \sqrt{\frac{\omega}{2\nu}}y
    \right)]
\end{equation}
To make sure the solution doesn't diverge at $y \rightarrow \infty$, we must take the negative root, giving:
\begin{equation}
    v(y, t) = v_0 \exp[-\sqrt{\frac{\omega}{2\nu}}y +i\left(
        \omega t - \sqrt{\frac{\omega}{2\nu}}y
    \right)]
\end{equation}
\subsection{Discussion}
The solution we got is a spatialy dacying wave with wave number $k = \sqrt{\frac{\omega}{2\nu}}$ and decay distance of $Y = \sqrt{\frac{2\nu}{\omega}}$. It fits our expectations of the frequency of the boundary conditions.
\newline
The original dimensionless variable from Stoke's first problem $\frac{y}{\sqrt{\nu t}}$ was replaced with $\pi_2 = \sqrt{\frac{\omega}{\nu}}y$. The fact that the system has two distinct constants allows it to create two distinct scales for distances and time, unlike Stoke's first problem, where there is a single constant and thus the distance and time scales are dependent. 

\newpage
\section{Venturi meter}
Since we are discussing steady flow under the influence of gravity, Bernoulli’s principle applies:
\begin{equation}
    \frac{1}{2}\rho v^2 + P + \rho h z = \text{const}
\end{equation}
\subsection{The pressure at the external points}
The external points $A', B', C'$ represent the boundary between the liquid and the air. This means that, by definition, the pressure at those points is the atmospheric pressure $p_0$.
\subsection{The velocity at the external points}
Since under the assumption there are only velocities in the $\boldsymbol{\hat{x}}$ direction, and since the 3 vertical tubes are too narrow for such currents, the total velocities there is zero. This can also be thought of from Bernoulli's principle at points $A', B', C'$, where the pressure is $p_0$, and equal to the external pressure, and so these points are at hydrostatic equillibrium.
\subsection{The pressure ODE in the main pipe}
To derive the ODE for the pressure in the main pipe, we look at Navier Stokes equation for incompressible fluids:
\begin{equation}
    \partial_t \boldsymbol{v} + (\boldsymbol{v}\cdot\boldsymbol{\nabla})\boldsymbol{v} = \boldsymbol{F} - \frac{1}{\rho}\boldsymbol{\nabla}p + \nu \nabla^2 \boldsymbol{v}
\end{equation}
Since it is an ideal flow, we can set $\nu = 0$ and since it is a steady flow we can set $\partial_t \boldsymbol{v} = 0$, giving us:
\begin{equation}
    (\boldsymbol{v}\cdot\boldsymbol{\nabla})\boldsymbol{v} = \boldsymbol{F} - \frac{1}{\rho}\boldsymbol{\nabla}p
\end{equation}
We also assume the velocity is constant and along the $\boldsymbol{\hat{x}}$ direction, which means:
\begin{equation}
    (\boldsymbol{v}\cdot\boldsymbol{\nabla})\boldsymbol{v} = 0
\end{equation}
So:
\begin{equation}
    \rho \boldsymbol{F} = \boldsymbol{\nabla}p
\end{equation}
Since the only force is gravity $-g \boldsymbol{\hat{z}}$, we get:
\begin{equation}
    \partial_x p = \partial_y p = 0 \rightarrow p = p(z)
\end{equation}
\begin{equation}
    \boxed{\frac{dp}{dz} + \rho g = 0}
\end{equation}
\subsection{Analysis}
\subsubsection{The narrow tubes}
In the narrow tubes, since there is hydrostatic equillibrium, the velocity cancels out and so just like in the previous section, the only remaining terms in the NS equaiton are the pressure and forces. Since again the only force working is in the z direction, we are left with the same pressure equation as in the last section.
\subsubsection{Height - pressure relation}
Using this equation, we can integrate over the height from the center of the large tube to the top of each narrow tube to get that for each of them:
\begin{equation}
    \Delta p = \int_0^{h} \rho g dz = \rho g h \rightarrow p_{center} = p_0 + \rho g h
\end{equation}
In other words:
\begin{equation}
    \boxed{
        p_A = p_0 + \rho g h_A
        \qquad
        p_B = p_0 + \rho g h_B
        \qquad
        p_C = p_0 + \rho g h_C
    }
\end{equation}
\subsubsection{Velocity - pressure relation}
At each point along the center of main pipe, we can apply the Bernoulli principle to get:
\begin{equation}
    \frac{1}{2}\rho v^2 + p = \text{const}
\end{equation}
Comparing points $A$ and $B$ gives:
\begin{equation}
    \boxed{p_A - p_B = \frac{1}{2}\rho\left(v_B^2 - v_A^2\right)}
\end{equation}
\subsubsection{Velocity - height relation}
Using the previous two equations, we can write:
\begin{equation}
    p_A - p_B = \rho g \left(h_A - h_B\right)
\end{equation}
Which gives the expected energy conservation equation:
\begin{equation}
    \boxed{g \left(h_A - h_B\right) = \frac{1}{2}\left(v_B^2 - v_A^2\right)}
\end{equation}
\subsubsection{$h_A$ and $h_C$ difference}
From symmetry considerations, if the both $A$ and $C$ have the same radius, there is no difference between them and so the pressure in these points would be the same, causing the same height difference. Let us show this using mass conservation.
\newline
If we take the boundary inside the large pipe and between points $A$ and $C$, the total flux must equal to change in the total mass of fluid inside:
\begin{equation}
    \partial_t \rho + \boldsymbol{\nabla}\cdot (\rho \boldsymbol{v}) = 0
\end{equation}
\begin{equation}
    \int_V \partial_t \rho d^3x + \int_V \boldsymbol{\nabla}\cdot (\rho \boldsymbol{v}) d^3x = 0
\end{equation}
\begin{equation}
    \partial_t m + \int_{\delta V} \rho \boldsymbol{v} \cdot d\boldsymbol{S} = 0
\end{equation}
Firstly, from ideal flow $\partial_t \rho = 0$. Secondly, we know that fluid cannot pass through the boundaries of the pipe, which means the only flux is in the $\boldsymbol{\hat{x}}$ direction with $v_x$. This gives:
\begin{equation}
    \int_{S} v_A d^2x = \int_{S} v_C d^2x \rightarrow v_A = v_C
\end{equation}
Pluggin this into the above equation gives:
\begin{equation}
    \boxed{h_A - h_C = \frac{1}{2g}(v_C^2 - v_A^2) = 0}
\end{equation}
\subsection{Variying density}
Now, we assume $\rho = \rho(z)$. We also assume $\boldsymbol{v} = \boldsymbol{\nabla}\phi$. Plugging this back into NS equation, the convection term changes to $(\boldsymbol{v}\cdot\boldsymbol{\nabla})\boldsymbol{v} = \frac{1}{2}\boldsymbol{\nabla}v^2$:
\begin{equation}
    \frac{1}{2}\boldsymbol{\nabla}v^2 = \boldsymbol{F} - \frac{1}{\rho}\boldsymbol{\nabla}p
\end{equation}
Assuming everything is consatnt in every section of the pipe, but changes with $z$, we get
\begin{equation}
    \frac{1}{2}\partial_z v_x^2 = g - \frac{1}{\rho}\partial_z p
\end{equation}
Integrating both sides along the $z$ direction from $z = 0$ to $z = h$ gives:
\begin{equation}
    \frac{1}{2}\left(v_x^2(h) - v_x^2(0)\right) = gh - \int_{p_0}^{p_h} \frac{dP}{\rho}
\end{equation}
Isolating $h$ and substracting $h_A - h_B$ gives:
\begin{equation}
    h_A - h_B = \frac{1}{2g}\left(v_A^2(h_A) - v_B^2(h_B)\right) + \int_{p_B}^{p_A} \frac{dP}{\rho g}
\end{equation}
Where $p_A, p_Bm v_A, v_B$ are the pressures and velocities at the center of the tubes at points $A$ and $B$, respectivly.
Finally, let us notice that since at points $h_A$ and $h_B$ the pressure and density are the same, we can use Bernoulli's principle between these two points and get and get that the velocities are also the same $v_A(h_A) = v_B(h_B)$. This results in
\begin{equation}
    h_A - h_B = \int_{p_B}^{p_A} \frac{dP}{\rho g}
\end{equation}
This integral goes along with a fluid parcel. Let us show that this parcel moves along the $x$ direction. Using the continuity equation with $\partial_t \rho = 0$, we get:
\begin{equation}
    \boldsymbol{\nabla}\cdot(\rho \boldsymbol{v}) = 0
\end{equation}
\begin{equation}
    \rho \boldsymbol{\nabla} \cdot \boldsymbol{v} + \boldsymbol{v}\cdot\boldsymbol{\nabla}\rho = 0
\end{equation}
Assuming $\rho = \rho(z)$ is not a constant, but the fluid is compressible $\boldsymbol{\nabla} \cdot \boldsymbol{v} = 0$, this gives:
\begin{equation}
    v_x \partial_x \rho + v_z \partial_z \rho = 0
\end{equation}
Since $\partial_x \rho = 0$, this requires $v_z = 0$. So, the fluid only moves horizontally. This means that the integral goes along a constant $z$ and thus a constant $\rho$. We can thus pull it out and get:
\begin{equation}
    \boxed{h_A - h_B = \frac{P_A - P_B}{\rho g}}
\end{equation}
Deriving the potential difference is just a matter of flipping the equation
\begin{equation}
    \boxed{P_A - P_B = \rho g\left(h_A - h_B\right)}
\end{equation}
\end{document}